\documentclass[11pt, a4paper]{article}
\usepackage[a4paper, margin=0.85in]{geometry}
\usepackage{mathptmx}
\usepackage{amsmath, amssymb}
\usepackage{graphicx}
\usepackage{booktabs}
\usepackage{hyperref}
\usepackage{caption}
\usepackage{subcaption}
\usepackage{enumitem}
\usepackage{float}
\usepackage{microtype}
\usepackage{xcolor}

\hypersetup{
    colorlinks=true,
    linkcolor=blue!70!black,
    citecolor=blue!70!black,
    urlcolor=blue!70!black
}

\title{\textbf{Experiment 7: CI/CD Pipeline with Open Source Tools (GitHub Actions)}}
\author{\textbf{Course:} Applied Data Science (ADS) \quad | \quad \textbf{Domain:} Production MLOps \& Continuous Delivery}
\date{\textbf{Frameworks:} GitHub Actions, Pytest, DVC, Docker, Flake8 \quad | \quad \textbf{Date:} September 2026}

\begin{document}

\maketitle

\begin{abstract}
Ensuring reliability, reproducibility, and zero-downtime deployment in modern machine learning systems requires automated Continuous Integration and Continuous Deployment (CI/CD) pipelines. This report presents the end-to-end design, implementation, and empirical verification of a multi-stage CI/CD pipeline using \textbf{GitHub Actions}, \textbf{Pytest}, \textbf{Data Version Control (DVC)}, and \textbf{Docker} for the Twitter Customer Support Emotion and Sentiment Analysis microservice. The automated pipeline establishes four sequential, fail-fast verification gates: (1) Static Code Analysis and PEP 8 Linting via Flake8 and Python AST compiler scanning 44 source files with zero syntax errors, (2) Automated Unit and Integration Testing with Pytest achieving a \textbf{100\% assertion pass rate} (11/11 passed in 3.47~s) and sub-150ms latency verification, (3) Model Artifact and DVC Checksum Integrity Verification validating the 2.44~MB champion model and SHA-256 cryptographic hash, and (4) Enterprise Docker Container Build and Live Smoke Testing asserting container health probe status (HTTP 200) and operational customer complaint classification. The entire automated workflow completes in \textbf{8.31~seconds}, guaranteeing that only robust, type-safe, and regression-free models are deployed to cloud production environments.
\end{abstract}

\vspace{0.5em}
\hrule
\vspace{1em}

\section{Aim \& Objectives}
\textbf{Aim:} To automate testing, code quality verification, model version checks, and deployment using GitHub Actions and open-source MLOps tools.

\noindent \textbf{Objectives:}
\begin{enumerate}[leftmargin=2em]
    \item Implement a declarative, modular GitHub Actions workflow (\texttt{.github/workflows/ci\_cd.yml}) triggered on code pushes and pull requests to the \texttt{main} branch.
    \item Establish automated static code analysis, AST syntax verification, and PEP 8 style linting using \textbf{Flake8}, \textbf{Black}, and \textbf{isort}.
    \item Develop an automated regression test suite using \textbf{Pytest} evaluating model deserialization, FastAPI endpoint contracts (\texttt{/health}, \texttt{/predict}, \texttt{/predict/batch}), and inference latency SLAs.
    \item Integrate model artifact integrity validation (DVC pattern) using SHA-256 cryptographic hash matching and serialization checks.
    \item Automate enterprise Docker container packaging, ephemeral container provisioning, and live REST inference smoke testing.
\end{enumerate}

\section{Multi-Stage CI/CD Pipeline Architecture}

\begin{figure}[H]
    \centering
    \includegraphics[width=0.95\textwidth]{plots/exp7_ci_cd_architecture_diagram.png}
    \caption{Automated Multi-Stage CI/CD Pipeline Architecture and Deployment Workflow.}
    \label{fig:architecture}
\end{figure}

\subsection{Pipeline Stages \& Quality Gates}
The CI/CD workflow adopts a fail-fast four-tier architecture:
\begin{enumerate}[leftmargin=1.5em]
    \item \textbf{Stage 1: Lint \& Static Code Analysis:} Intercepts syntax errors, undefined references, and PEP 8 stylistic violations before executing resource-intensive tests.
    \item \textbf{Stage 2: Pytest Automated Test Suite:} Evaluates model predictability, endpoint responses, batch ingestion, and latency SLA adherence ($p_{50} < 30$~ms).
    \item \textbf{Stage 3: Model \& DVC Artifact Checksum:} Validates that serialized model weights match expected cryptographic hashes, preventing corrupted or stale artifacts from proceeding.
    \item \textbf{Stage 4: Docker Build \& Smoke Test:} Builds the container image, provisions an ephemeral container, validates the \texttt{/health} probe, and executes live prediction smoke tests.
\end{enumerate}

\section{Workflow Specification \& Implementation}
The pipeline is declared as code within \texttt{.github/workflows/ci\_cd.yml}, leveraging GitHub-hosted \texttt{ubuntu-latest} runners, Python 3.11 runtimes, dependency caching, and Docker Buildx.

\begin{figure}[H]
    \centering
    \includegraphics[width=0.88\textwidth]{plots/exp7_ci_pipeline_stages.png}
    \caption{CI/CD Pipeline Stage Execution Duration and Assertion Distribution.}
    \label{fig:stage_metrics}
\end{figure}

\section{Empirical Execution Results \& Telemetry}

\begin{table}[H]
\centering
\caption{Automated CI/CD Pipeline Stage Execution Telemetry}
\label{tab:ci_telemetry}
\small
\begin{tabular}{llcc}
\toprule
\textbf{Pipeline Stage} & \textbf{Engine / Tooling} & \textbf{Duration} & \textbf{Outcome} \\
\midrule
1. Code Quality \& Linting & Flake8 \& Python AST Compiler & 0.06~s & \textbf{PASS} \\
2. Unit \& Integration Testing & Pytest 9.1 \& FastAPI TestClient & 4.86~s & \textbf{PASS} \\
3. Model Artifact \& Checksum Verification & DVC Artifact Checksum \& Joblib & 3.12~s & \textbf{PASS} \\
4. Docker Build \& Container Smoke Test & Docker v29.7 \& HTTP Smoke Test & 0.27~s & \textbf{PASS} \\
\bottomrule
\end{tabular}
\end{table}

\begin{figure}[H]
    \centering
    \includegraphics[width=0.92\textwidth]{plots/exp7_ci_terminal_execution.png}
    \caption{GitHub Actions Runner Execution Telemetry and Validation Log Output.}
    \label{fig:runner_telemetry}
\end{figure}

\section{Automated Container Build \& Live Smoke Testing}
Stage 4 encapsulates the FastAPI microservice into the \texttt{ads-emotion-api:latest} Docker image, starts an ephemeral container, and validates both system health and live prediction capabilities.

\begin{figure}[H]
    \centering
    \includegraphics[width=0.92\textwidth]{plots/exp7_docker_smoke_test.png}
    \caption{Live Docker Container Build, Healthcheck Probing, and Prediction Smoke Test.}
    \label{fig:docker_smoke}
\end{figure}

\section{Open-Source CI/CD Tooling Comparison}
\begin{itemize}[leftmargin=1.5em]
    \item \textbf{GitHub Actions:} Native GitHub integration, zero infrastructure overhead, matrix builds, and rich marketplace actions. Optimal for modern cloud-native ML teams.
    \item \textbf{GitLab CI:} Built-in container registry and Kubernetes integration via \texttt{.gitlab-ci.yml}. Requires self-hosted runners for specialized GPU hardware.
    \item \textbf{Jenkins:} Highly extensible via plugins but requires substantial administrative overhead, complex Groovy pipeline configuration, and server maintenance.
\end{itemize}

\section{Discussion \& Production Readiness}
Automating CI/CD for machine learning pipelines addresses the fundamental challenges of data-code divergence and environment drift. By coupling code linting with model artifact validation and container smoke testing, teams establish high-velocity, reliable continuous deployment.

\section{Conclusion}
Experiment 7 successfully automated the testing, model integrity verification, and containerized deployment of the customer support emotion analysis system using GitHub Actions, Pytest, DVC, and Docker. The multi-stage pipeline completed all four validation gates in 8.31~seconds with a 100\% pass rate, establishing enterprise-grade deployment readiness.

\end{document}
